\documentclass[letterpaper,11pt]{article}
\usepackage[margin=0.9in]{geometry}
\usepackage{hyperref}
\pagestyle{empty}
\setlength{\parindent}{0pt}
\setlength{\parskip}{6pt}

\begin{document}

{\Huge \scshape Vedaang Chopra}\\[2pt]
\small vedaangchopra@gatech.edu \;|\; +1 (404) 740-9905 \;|\; Atlanta, GA \;|\;
\href{https://linkedin.com/in/vedaang-chopra}{LinkedIn} \;|\;
\href{https://github.com/Vedaang-Chopra}{GitHub}

\vspace{10pt}
\hrule
\vspace{14pt}

Dear Scale AI Hiring Team,

I am applying for the Frontier Agents Engineer (Applied AI) position. My work sits exactly at the intersection this role demands: building agentic AI systems that survive contact with production, and building the evaluation harnesses that prove they work.

\textbf{Production agents, not prototypes.} At Fortinet's AIOps R\&D, I led an agentic RAG diagnostics system from hackathon prototype to production deployment: an LLM that plans multi-step tool calls over live network telemetry, invokes APIs, and verifies hypotheses through structured function calling to produce autonomous root-cause analysis -- cutting mean resolution time by ${\sim}70\%$. I scaled OpenSearch pipelines from 50 to 2,000 events/sec along the way, so I have felt the reliability, latency, and observability constraints your enterprise customers operate under.

\textbf{Multi-agent architecture with real orchestration.} My research spans three systems: ATHENA, a five-agent supervisor-worker system with reflection-based quality critique and dynamic plan modification; SHASTRA (in progress), a framework for retrieving and adapting prior agent workflows for new long-horizon tasks under cost/latency budgets; and a closed-loop CAD code-generation pipeline where an LLM's output is executed, verified geometrically, and repaired iteratively -- agents grounded by verifiable execution feedback, not single-pass generation.

\textbf{Evaluation as a first-class discipline.} I maintain a 202-test evaluation harness with golden datasets, regression suites, and LLM-as-a-Judge visual scoring that gates every change to my current lab's 16-module pipeline. I have run controlled ablations -- removing my ATHENA system's research agent cut fact coverage by 27\% in human-rated evaluation (n=5 raters) -- because ``own the full experimentation lifecycle'' is not a slogan; it is how I already work.

What draws me to Scale is the portfolio breadth -- solving many different AI problems across industries rather than one narrow slice -- and the mandate to pair frontier-model fluency with deterministic engineering. That pairing matches my track record across Python, Go, distributed systems (Azure-deployed), and modern agent tooling.

I would welcome the chance to discuss how my experience building trustworthy agentic systems can contribute to Scale's enterprise deployments.

\vspace{8pt}
Sincerely,\\
Vedaang Chopra

\end{document}
